\guidechapter{Small error must stay small after rescaling}{A numerical budget you can audit by hand.}

Imagine an approximate computation whose current absolute error is at most
$0.01$. Passing it through an exact contraction does not enlarge that error.
If the next approximation adds at most $0.02$, the new error bound is at
most $0.03$. This is a worst-case sum: we do not assume errors cancel or
behave randomly.

Now restore a scale factor of eight. The error bound becomes
$8\cdot0.03=0.24$, which is safely below one half. If the scale factor
were 32 instead, the bound would become $0.96$. That weaker guarantee does
not justify rounding to the right integer.

\explfigure{20-error-budget}{A small internal error is not the final error.}{
The last two bars are alternative rescalings of the same $0.03$ bound,
not successive operations. Crossing $1/2$ means this bound is insufficient
to guarantee exact rounding; it does not prove the actual answer is wrong.}

The source proof has larger and more elaborate normalizations, but the
principle is the same. It budgets local errors, bounds intermediate sizes,
and tracks the multiplication of error when scales are undone. It also
needs enough bits for intermediate exact dyadic values, even when the
completed transform is contractive.

\paragraph{Why strict inequality matters.}
If the true integer is 12, a result strictly between $11.5$ and $12.5$
rounds uniquely to 12. At exactly $12.5$, the distance to 12 and 13 is the
same, so an error allowance of ``at most one half'' cannot determine which
integer was intended. The paper proves a strict final bound for every
coefficient and then performs exact carry propagation.

\takeaway{Approximation is compatible with an exact answer when the
fully rescaled error is proved smaller than half the gap between integers.}
